% !TEX program = pdflatex
% ============================================================
% Report esame Corso LPSMT — Progetto PagUp
% Compilazione: pdflatex report.tex (eseguire due volte per i riferimenti)
% ============================================================
\documentclass[11pt,a4paper]{article}

\usepackage[utf8]{inputenc}
\usepackage[T1]{fontenc}
\usepackage[italian]{babel}
\usepackage[a4paper,margin=2.5cm]{geometry}
\usepackage{lmodern}
\usepackage{microtype}
\usepackage{xcolor}
\usepackage{graphicx}
\usepackage{booktabs}
\usepackage{tabularx}
\usepackage{enumitem}
\usepackage{fancyhdr}
\usepackage{titlesec}
\usepackage{caption}
\usepackage{float}
\usepackage{tikz}
\usetikzlibrary{shapes.geometric, arrows.meta, positioning, fit, backgrounds, calc}
\usepackage[hidelinks]{hyperref}

% ---------- Palette (i colori del design system dell'app) ----------
\definecolor{pagprimary}{HTML}{003A80}
\definecolor{paggreen}{HTML}{145A32}
\definecolor{paggrey}{HTML}{303030}
\definecolor{paglight}{HTML}{EBEBEB}
\definecolor{pagcode}{HTML}{F5F5F5}

% ---------- Header / footer ----------
\pagestyle{fancy}
\fancyhf{}
\fancyhead[L]{\small\textsc{PagUp}}
\fancyhead[R]{\small Report esame LPSMT}
\fancyfoot[C]{\thepage}
\renewcommand{\headrulewidth}{0.4pt}

% ---------- Titoli ----------
\titleformat{\section}{\Large\bfseries\color{pagprimary}}{\thesection.}{0.6em}{}
\titleformat{\subsection}{\large\bfseries\color{paggrey}}{\thesubsection}{0.6em}{}

% ---------- Stili TikZ ----------
\tikzset{
  node distance=1.1cm and 1.3cm,
  box/.style={rectangle, rounded corners=3pt, draw=pagprimary, thick,
    fill=pagprimary!6, align=center, minimum height=1cm, inner sep=6pt,
    text width=2.8cm, font=\small},
  store/.style={cylinder, shape border rotate=90, aspect=0.3, draw=paggreen,
    thick, fill=paggreen!8, align=center, minimum height=1.4cm, text width=2.3cm, font=\small},
  ext/.style={rectangle, rounded corners=3pt, draw=paggrey, thick,
    fill=paglight, align=center, minimum height=0.9cm, text width=2.6cm, font=\small},
  flow/.style={-{Stealth[length=2.5mm]}, thick, draw=paggrey},
  bidir/.style={{Stealth[length=2.5mm]}-{Stealth[length=2.5mm]}, thick, draw=pagprimary},
}

% ============================================================
\begin{document}

% ---------- Frontespizio ----------
\begin{titlepage}
\centering
\vspace*{2cm}
{\Huge\bfseries\color{pagprimary} PagUp\par}
\vspace{0.4cm}
{\large Autonomia e inclusione nel pagamento in contanti\par}
\vspace{0.2cm}
{\normalsize Wallet educativo e accessibile per persone con disabilit\`a cognitiva\par}
\vspace{2cm}
{\large\bfseries Report di progetto\par}
{\large Corso LPSMT --- Laboratorio di Programmazione per Sistemi Mobili e Tablet\par}
\vspace{2.5cm}
{\Large Paolo Sarcletti \quad\textperiodcentered\quad Alessandro Turri\par}
\vspace{2cm}
\begin{tabular}{rl}
\textbf{Tecnologie:} & Expo (React Native) + TypeScript, Supabase \\[4pt]
\textbf{Piattaforme:} & Android, iOS \\[4pt]
\textbf{Anno accademico:} & 2025/2026 \\
\end{tabular}
\vfill
\end{titlepage}

\tableofcontents
\newpage

% ============================================================
\section{Introduzione}

\subsection{Contesto dell'applicazione e NEED identificato}

Pagare in contanti alla cassa \`e un gesto che la maggior parte delle persone compie senza
pensarci. In realt\`a \`e un compito cognitivo composito: bisogna riconoscere i tagli di monete e
banconote, comporre un importo a partire dai pezzi disponibili nel portafoglio, valutare se la
cifra posseduta \`e sufficiente e infine controllare il resto. Per una persona con
\textbf{disabilit\`a cognitiva} (disabilit\`a intellettiva, sindrome di Down, disturbi
dell'apprendimento) ognuno di questi passaggi pu\`o trasformarsi in un ostacolo, e l'ostacolo si
presenta nel momento peggiore: in pubblico, con la fila che attende. Il risultato tipico \`e che
l'esperienza dell'acquisto viene evitata o delegata per intero al caregiver, con un costo che non
\`e solo pratico ma anche psicologico, perch\'e la dipendenza per ogni acquisto erode il senso di
autoefficacia della persona.

Dal lato opposto della relazione c'\`e il tutor --- un genitore, un educatore, un operatore di
centro diurno --- che vorrebbe accompagnare gradualmente la persona verso l'autonomia, ma che oggi
pu\`o farlo solo stando fisicamente accanto a lei. Nel momento in cui la persona esce da sola, il
tutor perde ogni visibilit\`a e ogni possibilit\`a di intervento.

Il NEED che abbiamo identificato nasce dall'incrocio di questi due bisogni, finora mai coperti
da un'unica soluzione:

\begin{enumerate}[leftmargin=1.4cm]
  \item \textbf{per lo studente}: uno strumento che guidi \emph{passo-passo} l'atto del pagamento
  reale, riduca il carico cognitivo a una sola decisione per volta, dia un riscontro immediato e
  multisensoriale (visivo ad alto contrasto, aptico, testuale semplificato) e permetta di
  allenarsi senza rischi prima di affrontare la cassa vera;
  \item \textbf{per il tutor}: la possibilit\`a di configurare il portafoglio digitale, calibrare
  la difficolt\`a dell'esercizio, osservare in tempo reale come procede e ricevere subito una
  richiesta di aiuto (SOS) georeferenziata quando la persona \`e in difficolt\`a fuori casa.
\end{enumerate}

In una frase: \emph{trasformare un momento di stress in un'azione guidata e sicura}, rendendo la
persona il pi\`u autonoma possibile senza togliere al tutor la rete di protezione.

\subsection{Obiettivi generali del progetto}

PagUp \`e un'applicazione mobile cross-platform (Android/iOS) con due ruoli distinti --- studente
e tutor --- costruiti sopra lo stesso backend. Gli obiettivi che ci siamo posti:

\begin{itemize}[leftmargin=1.4cm]
  \item \textbf{Accessibilit\`a come requisito di prima classe}, non come rifinitura: l'area
  studente rispetta il livello \emph{AAA delle WCAG~2.1} (contrasto $\geq 7{:}1$ su ogni coppia
  testo/sfondo), con touch target da 64\,dp (oltre il minimo di 48\,dp), tipografia mai sotto i
  18\,px, etichette di accessibilit\`a e feedback aptico su ogni azione.
  \item \textbf{Wizard di pagamento guidato} in quattro passi (importo $\rightarrow$ conferma
  $\rightarrow$ quali pezzi consegnare $\rightarrow$ esito e resto), con un algoritmo che calcola
  automaticamente la combinazione di monete e banconote da usare.
  \item \textbf{Input multimodale} dell'importo: tastierino numerico, lettura OCR del totale dallo
  scontrino tramite fotocamera, dettatura vocale in italiano. Chi non riesce a digitare pu\`o
  fotografare o parlare.
  \item \textbf{Modalit\`a allenamento}: stesso gesto del pagamento, cifre generate casualmente,
  nessuna conseguenza sul saldo reale.
  \item \textbf{Supervisione remota in tempo reale}: dashboard tutor con saldo, statistiche,
  cronologia delle attivit\`a, editor del portafoglio e scelta della modalit\`a di pagamento,
  sincronizzata via Supabase Realtime senza polling.
  \item \textbf{Rete di protezione}: pulsante SOS sempre visibile durante il pagamento, con invio
  di posizione GPS e notifica push al tutor; impostazioni dello studente protette da PIN del
  tutor; isolamento dei dati garantito da Row Level Security nel database.
\end{itemize}

% ============================================================
\section{Identificazione del segmento utente}

\subsection{Utenti target}

L'applicazione \`e progettata attorno a una relazione uno-a-molti tra tutor e studenti. I due
profili condividono il backend ma usano interfacce, temi grafici e flussi completamente separati.

\begin{table}[H]
\centering
\renewcommand{\arraystretch}{1.35}
\begin{tabularx}{\textwidth}{@{}p{3.2cm}XX@{}}
\toprule
 & \textbf{Studente (utente primario)} & \textbf{Tutor (utente secondario)} \\
\midrule
\textbf{Chi \`e} & Persona con disabilit\`a cognitiva lieve o media in percorso verso
l'autonomia; in prospettiva anche bambini in et\`a scolare che imparano a gestire il denaro
& Genitore, educatore, operatore di centro diurno, insegnante di sostegno \\
\textbf{Obiettivo} & Pagare e gestire il denaro in autonomia & Accompagnare, configurare,
monitorare, intervenire \\
\textbf{Interfaccia} & Tema \texttt{studentTheme}: contrasto AAA, touch target 64\,dp, testo
$\geq$18\,px, una sola azione per schermata & Tema \texttt{tutorTheme}: dashboard informativa a
quattro schede (home, statistiche, studenti, wallet) \\
\textbf{Schermate} & Home, pagamento guidato, allenamento & Dashboard, impostazioni (associazione
via QR, gestione PIN) \\
\bottomrule
\end{tabularx}
\caption{I due segmenti utente di PagUp e le rispettive caratteristiche.}
\end{table}

\subsection{Perch\'e proprio loro: la lacuna da colmare}

La scelta del segmento risponde a lacune precise del panorama attuale. Le app finanziarie
generaliste (home banking, wallet digitali, gestori di spese) sono progettate per utenti
neurotipici: presuppongono lettura fluente, astrazione numerica e tollerano interfacce dense; per
il nostro profilo target risultano inaccessibili quando non ansiogene. Le app didattiche sul
denaro, all'estremo opposto, trattano il denaro come concetto astratto --- quiz e conteggi a
schermo --- ma non accompagnano la persona nel momento in cui il pagamento avviene davvero, e non
prevedono alcuna figura di supervisione.

Manca cio\`e una soluzione che tenga insieme tre cose: un'accessibilit\`a di livello verificabile
(non ``caratteri grandi'' generici, ma contrasti e target misurati), un \emph{ponte con il denaro
fisico} --- l'app che dice concretamente quali pezzi prendere dal portafoglio vero, mostrandone la
fotografia --- e una supervisione remota del caregiver. \`E esattamente lo spazio che PagUp
occupa. Si tratta di un'utenza fragile e sottoservita, per la quale il costo dell'errore
(vergogna, frustrazione, perdita di denaro) \`e alto e ogni dettaglio di accessibilit\`a decide se
l'app \`e usabile oppure no.

% ============================================================
\section{Stato dell'arte}

\subsection{Applicazioni simili presenti negli store}

Il panorama esistente si pu\`o suddividere in quattro famiglie, nessuna delle quali copre per
intero il bisogno individuato.

\begin{enumerate}[leftmargin=1.4cm]
  \item \textbf{App di riconoscimento del denaro per disabilit\`a visiva} --- ad esempio
  \emph{Cash Reader} o il canale valuta di \emph{Seeing AI} (Microsoft). Identificano una
  banconota inquadrata con la fotocamera e la annunciano vocalmente. Sono tecnicamente mature, ma
  nascono per persone cieche o ipovedenti: riconoscono \emph{un} pezzo, non aiutano a
  \emph{comporre} un importo n\'e a decidere cosa consegnare alla cassa.
  \item \textbf{App didattiche sul denaro} --- esercizi di conteggio di monete ed euro pensati
  per la scuola primaria o per la didattica speciale. Insegnano il riconoscimento dei tagli
  tramite quiz, ma l'esercizio resta scollegato dall'atto reale del pagamento, l'interazione \`e
  spesso datata, la sincronizzazione cloud assente e la conformit\`a a standard di accessibilit\`a
  verificabili pressoch\'e mai dichiarata.
  \item \textbf{App di banking e wallet generaliste} --- home banking, Google Wallet, Satispay e
  simili. Maturit\`a tecnica molto elevata, ma fuori target: interfacce dense, gergo finanziario,
  denaro esclusivamente virtuale, nessuna attenzione all'accessibilit\`a cognitiva.
  \item \textbf{App di assistenza e comunicazione per la disabilit\`a} --- strumenti di CAA
  (comunicazione aumentativa alternativa), agende visive, app di sicurezza con pulsante di
  emergenza. Coprono aspetti adiacenti (comunicazione, allerta al caregiver) ma non affrontano il
  dominio della gestione del denaro.
\end{enumerate}

\subsection{Livello di maturazione}

Le famiglie 1 e 3 presentano un livello di maturazione elevato, ma servono target diversi dal
nostro (disabilit\`a visiva la prima, popolazione generale la seconda). La famiglia 2, che \`e la
pi\`u vicina per intento educativo, \`e anche la pi\`u debole per qualit\`a: molte app risultano
incomplete (coprono il riconoscimento ma non il pagamento), abbandonate o prive di qualunque
componente di supervisione. La combinazione \emph{didattica del denaro + accompagnamento nel
pagamento reale + supervisione del caregiver in tempo reale} \`e, per quanto abbiamo potuto
verificare negli store, semplicemente assente.

\subsection{Posizionamento di PagUp}

\begin{table}[H]
\centering
\renewcommand{\arraystretch}{1.3}
\begin{tabularx}{\textwidth}{@{}p{4.6cm}cccc@{}}
\toprule
\textbf{Caratteristica} & \textbf{Didattiche} & \textbf{Banking} & \textbf{CAA/SOS} & \textbf{PagUp} \\
\midrule
Accessibilit\`a WCAG AAA          & parziale & no & parziale & \textbf{s\`i} \\
Wizard di pagamento guidato       & no       & no & no       & \textbf{s\`i} \\
Ponte col denaro fisico reale     & no       & no & no       & \textbf{s\`i} \\
Modalit\`a allenamento sicura     & s\`i     & no & no       & \textbf{s\`i} \\
Supervisione tutor in tempo reale & no       & no & no       & \textbf{s\`i} \\
SOS con geolocalizzazione         & no       & no & s\`i     & \textbf{s\`i} \\
Input OCR + voce                  & no       & raro & no     & \textbf{s\`i} \\
\bottomrule
\end{tabularx}
\caption{Posizionamento di PagUp rispetto alle famiglie di applicazioni esistenti.}
\end{table}

PagUp non compete con il banking sulla gestione finanziaria n\'e con le app didattiche sul puro
esercizio: si propone come \emph{strumento ponte verso l'autonomia reale}, ed \`e l'unico prodotto
della tabella a coprire l'intera filiera che va dall'allenamento protetto al pagamento vero
supervisionato a distanza.

% ============================================================
\section{Architettura}

\subsection{Visione d'insieme}

PagUp adotta un'architettura \textbf{client--serverless} di tipo BaaS (\emph{Backend as a
Service}). Il client \`e un'unica applicazione Expo / React Native scritta in TypeScript che, in
base al ruolo dell'utente autenticato, instrada verso l'esperienza studente o quella tutor. Il
backend \`e interamente delegato a \textbf{Supabase}, che sopra PostgreSQL fornisce
autenticazione, persistenza, canali Realtime e sicurezza a livello di riga (Row Level Security).
Non esiste un server applicativo intermedio: la logica di dominio vive nel client, isolata in
funzioni pure e testata con Jest, mentre le regole di accesso ai dati vivono nel database sotto
forma di policy dichiarative.

\begin{table}[H]
\centering
\renewcommand{\arraystretch}{1.25}
\begin{tabularx}{\textwidth}{@{}p{3.6cm}p{4.4cm}X@{}}
\toprule
\textbf{Livello} & \textbf{Tecnologia} & \textbf{Ruolo} \\
\midrule
Frontend & Expo SDK 54, React Native 0.81, React 19, TypeScript & UI, navigazione, logica di dominio \\
Stato locale & Zustand 5 + persistenza AsyncStorage & Portafoglio dello studente, sottoscrizione Realtime \\
Navigazione & React Navigation (Stack) & Routing per ruolo, tasto indietro su ogni schermata non-home \\
Backend & Supabase (PostgreSQL) & Auth, dati, Realtime, RLS \\
Notifiche & Expo Notifications + Expo Push API & Avvisi push al tutor (pagamenti, SOS) \\
Servizi nativi & ML Kit (OCR), Speech Recognition, Location, Camera, Haptics & Input multimodale, SOS, feedback tattile \\
\bottomrule
\end{tabularx}
\caption{Stack tecnologico di PagUp.}
\end{table}

\subsection{Diagramma dell'architettura generale}

\begin{figure}[H]
\centering
\begin{tikzpicture}
  % --- Client layer ---
  \node[box] (student) {App\\Studente};
  \node[box, right=2.6cm of student] (tutor) {App\\Tutor};
  \node[box, fill=paggreen!8, draw=paggreen, below=0.9cm of student] (zustand) {Zustand\\+ AsyncStorage};

  % --- External native services (left) ---
  \node[ext, left=1.8cm of student, yshift=1.2cm] (ocr) {ML Kit OCR};
  \node[ext, left=1.8cm of student, yshift=-0.1cm] (voice) {Speech\\Recognition};
  \node[ext, left=1.8cm of student, yshift=-1.4cm] (geo) {Location\\(SOS)};

  % --- Supabase core ---
  \node[store, below=2.6cm of $(student)!0.5!(tutor)$] (db) {Supabase\\PostgreSQL};
  \node[ext, draw=paggreen, fill=paggreen!6, left=1.1cm of db] (auth) {Supabase\\Auth};
  \node[ext, draw=paggreen, fill=paggreen!6, right=1.1cm of db] (rt) {Supabase\\Realtime};

  % --- Push ---
  \node[ext, below=1.0cm of tutor, xshift=2.2cm] (push) {Expo\\Push API};

  % --- Edges ---
  \draw[flow] (ocr) -- (student);
  \draw[flow] (voice) -- (student);
  \draw[flow] (geo) -- (student);
  \draw[bidir] (student) -- (zustand);
  \draw[bidir] (student) -- (db);
  \draw[bidir] (tutor) -- (db);
  \draw[bidir] (db) -- (auth);
  \draw[bidir] (db) -- (rt);
  \draw[flow] (rt) to[bend left=12] (student);
  \draw[flow] (rt) to[bend right=12] (tutor);
  \draw[flow] (student.south east) to[bend right=18] (push);
  \draw[flow] (push) -- (tutor);

  \begin{scope}[on background layer]
    \node[draw=pagprimary, dashed, rounded corners, fit=(student)(tutor)(zustand), inner sep=10pt, label={[pagprimary]above:\small Client Expo / React Native}] {};
    \node[draw=paggreen, dashed, rounded corners, fit=(db)(auth)(rt), inner sep=10pt, label={[paggreen]below:\small Backend Supabase (BaaS)}] {};
  \end{scope}
\end{tikzpicture}
\caption{Architettura generale: un unico client a doppio ruolo, servizi nativi del dispositivo,
backend Supabase e canale di notifiche push.}
\end{figure}

L'associazione tra un tutor e uno studente avviene di persona e senza scambio di credenziali: lo
studente mostra un QR code generato dall'app (contenente il proprio identificativo), il tutor lo
scansiona dalla schermata impostazioni e la riga corrispondente viene creata nella tabella ponte.
Da quel momento le policy RLS aprono al tutor la visibilit\`a sui dati di quello studente, e
soltanto su quelli.

\subsection{Modello dati e sicurezza (Row Level Security)}

Il database \`e composto da cinque tabelle. La sicurezza non \`e affidata al client ma
\emph{dichiarata nel database}: ogni utente autenticato pu\`o leggere e scrivere solo le righe che
le policy gli concedono, in funzione di \texttt{auth.uid()} e della tabella ponte
\texttt{tutor\_students}.

\begin{table}[H]
\centering
\renewcommand{\arraystretch}{1.25}
\begin{tabularx}{\textwidth}{@{}p{3.0cm}X@{}}
\toprule
\textbf{Tabella} & \textbf{Contenuto e regole principali} \\
\midrule
\texttt{profiles} & Estende \texttt{auth.users}: ruolo, username, nome completo, \texttt{tutor\_pin}. Creata da trigger al signup. Lo studente legge i profili dei propri tutor; il tutor quelli dei propri studenti. \\
\texttt{wallets} & Un portafoglio per studente; campo \texttt{items} JSONB con l'array di \texttt{MoneyItem} (\texttt{id}, \texttt{value}, \texttt{type}, \texttt{imageUri}). Lo studente aggiorna il proprio; il tutor legge e scrive quelli dei propri studenti. \\
\texttt{tutor\_students} & Tabella ponte della relazione tutor--studente; contiene la modalit\`a di pagamento assegnata (\texttt{payment\_mode}: \texttt{exact}/\texttt{fast}). Popolata tramite la scansione del QR. \\
\texttt{activity\_logs} & Cronologia eventi: \texttt{payment}, \texttt{sos}, \texttt{wallet\_adjustment}. Per i pagamenti registra importo richiesto, importo consegnato e modalit\`a usata; per gli SOS la schermata di provenienza, l'importo eventualmente inserito e la posizione GPS (in \texttt{metadata} JSONB). \\
\texttt{push\_tokens} & Token Expo Push per utente e piattaforma, usati per recapitare gli avvisi al tutor. \\
\bottomrule
\end{tabularx}
\caption{Schema dati e regole di accesso (RLS) di PagUp.}
\end{table}


\subsection{Sincronizzazione in tempo reale (broadcast)}

Il punto qualificante dell'architettura \`e la sincronizzazione bidirezionale \emph{senza polling}
tramite Supabase Realtime, che ritrasmette ai client sottoscritti gli eventi
\texttt{postgres\_changes} del database. Quando il tutor ricarica il portafoglio di uno studente,
la riga \texttt{wallets} cambia, l'evento \texttt{UPDATE} raggiunge il canale a cui il dispositivo
dello studente \`e sottoscritto, l'inventario locale si aggiorna e parte un feedback aptico che
avvisa la persona senza bisogno di leggere nulla. Simmetricamente la dashboard del tutor \`e
sottoscritta agli eventi su \texttt{wallets} e \texttt{activity\_logs}: un pagamento o un SOS
compaiono sullo schermo del tutor nel momento in cui avvengono. Lo stesso meccanismo sblocca in
automatico la home dello studente nell'istante in cui il tutor ne scansiona il QR.

\begin{figure}[H]
\centering
\begin{tikzpicture}[node distance=1.2cm and 2.0cm]
  \node[box] (tutorapp) {Tutor\\ricarica wallet};
  \node[store, right=of tutorapp] (db) {Supabase\\\texttt{wallets}};
  \node[box, right=of db] (studentapp) {App Studente\\(canale \texttt{wallet:uid})};
  \node[box, below=1.0cm of studentapp] (zs) {Zustand:\\inventario aggiornato};
  \node[ext, below=1.0cm of zs] (hap) {Feedback aptico\\Success};

  \draw[flow] (tutorapp) -- node[above,font=\scriptsize]{UPDATE} (db);
  \draw[flow] (db) -- node[above,font=\scriptsize]{postgres\_changes} (studentapp);
  \draw[flow] (studentapp) -- (zs);
  \draw[flow] (zs) -- (hap);
\end{tikzpicture}
\caption{Flusso di broadcast Realtime: la modifica del tutor si propaga al dispositivo dello
studente senza polling.}
\end{figure}


\subsection{Logica di dominio: l'algoritmo di pagamento}

Il cuore funzionale \`e il calcolo di \emph{quali pezzi} consegnare per coprire un importo. Tutti
i conti avvengono in \textbf{centesimi interi}, per evitare gli errori di arrotondamento del
floating point, e l'importo obiettivo viene arrotondato per eccesso al multiplo di 10 centesimi:
i tagli da 1, 2 e 5 centesimi sono esclusi sia dall'algoritmo sia dall'editor del tutor, perch\'e
maneggiarli aggiunge complessit\`a senza valore didattico. Le modalit\`a sono due, e le sceglie il
tutor per ciascuno studente:

\begin{itemize}[leftmargin=1.4cm]
  \item \textbf{Precisa (exact)}: cerca prima una combinazione \emph{esatta} dell'importo
  (variante di subset-sum con memoizzazione delle somme raggiungibili in centesimi); se non
  esiste, ripiega sulla pi\`u piccola combinazione che copre l'importo, minimizzando il resto.
  \item \textbf{Veloce (fast)}: privilegia la semplicit\`a cognitiva scegliendo il \emph{singolo}
  taglio pi\`u piccolo che da solo copre l'importo (per pagare 3{,}67~euro l'app propone la
  banconota da 5~euro). \`E il gradino iniziale del percorso: prima si impara ``dai una
  banconota abbastanza grande e aspetta il resto'', poi si passa alla composizione precisa.
\end{itemize}


La logica di dominio \`e isolata in funzioni pure (\texttt{src/utils/paymentLogic.ts},
\texttt{smartExtract.ts}) e coperta da test unitari Jest (combinazione esatta, fallback, modalit\`a
veloce, esclusione delle monete piccole, casi di insufficienza), secondo la regola di progetto
``hooks e utility (logica) separati dai componenti (UI)''.

\subsection{Flusso del wizard di pagamento}

\begin{figure}[H]
\centering
\begin{tikzpicture}[node distance=0.8cm and 1.5cm]
  \node[box] (a) {1. Importo\\(tastiera/OCR/voce)};
  \node[box, right=of a] (b) {2. Conferma\\importo};
  \node[box, right=of b] (c) {3. Quali pezzi\\consegnare};
  \node[box, right=of c] (d) {4. Esito\\+ resto};
  \node[ext, below=1.0cm of b, xshift=1.5cm] (rec) {recordPayment\\$\rightarrow$ log + push};
  \node[ext, below=1.0cm of a] (sos) {SOS\\(posizione + push)};

  \draw[flow] (a) -- (b);
  \draw[flow] (b) -- (c);
  \draw[flow] (c) -- (d);
  \draw[flow] (c) -- (rec);
  \draw[flow] (a.south) -- (sos);
\end{tikzpicture}
\caption{Macchina a stati del PaymentWizard: ogni passo \`e una schermata a s\'e, con pulsante SOS
sempre disponibile e registrazione dell'evento alla conferma del pagamento.}
\end{figure}

Ogni passo occupa l'intero schermo e chiede una sola decisione; si avanza e si torna indietro con
animazioni direzionali coerenti (avanti = da destra, indietro = da sinistra) che aiutano a
mantenere l'orientamento nel flusso. Nel terzo passo l'app mostra le \emph{fotografie} dei pezzi
da prendere dal portafoglio reale. Alla conferma, \texttt{processRealPayment} sottrae i pezzi
dall'inventario, persiste il wallet su Supabase, registra un \texttt{activity\_log} per ciascun
tutor associato --- con importo richiesto, importo consegnato e modalit\`a usata --- e invia la
notifica push. La schermata finale, a sfondo verde pieno, comunica l'esito e il resto atteso con
un linguaggio ridotto all'essenziale.

\subsection{L'accessibilit\`a come vincolo architetturale}

L'accessibilit\`a non \`e uno strato cosmetico ma una scelta d'architettura: l'app definisce due
\emph{design system} distinti come moduli TypeScript tipizzati (\texttt{src/theme/student.ts} e
\texttt{tutor.ts}), e nessuna schermata usa colori o dimensioni fuori dai token del proprio tema.
Il tema studente dichiara nel codice, accanto a ogni colore, il rapporto di contrasto verificato
con la formula di luminanza relativa delle WCAG (dal 7{,}0:1 minimo del testo disabilitato al 21:1
del nero su bianco); fissa il touch target a 64\,dp con distanza minima di 16\,dp tra bersagli
adiacenti, e il corpo del testo a 18\,px con interlinea 1{,}6. Ogni elemento interattivo espone
\texttt{accessibilityLabel} e \texttt{accessibilityHint} in italiano semplice, e ogni azione
produce un feedback aptico differenziato (selezione, conferma, successo, avviso, errore), cos\`i
che l'esito di un gesto sia percepibile anche senza leggere. Il tema tutor, rivolto a un utente
neurotipico, adotta il livello AA e una densit\`a informativa maggiore: la separazione dei due
temi permette di ottimizzare ciascuna esperienza senza compromessi reciproci.

% ============================================================
\section{Analisi critica dei possibili limiti}

Pur coprendo il bisogno individuato, l'applicazione ha limiti noti. Li analizziamo nel dettaglio
insieme alle direzioni di miglioramento, distinguendo ci\`o che \`e una scelta di scope da ci\`o
che \`e un debito tecnico.

\subsection{Limiti funzionali e di dominio}
\begin{itemize}[leftmargin=1.4cm]
  \item \textbf{Il resto \`e solo indicativo.} L'app calcola quanto resto attendersi ma non pu\`o
  verificare che il cassiere lo restituisca correttamente: l'anello finale dell'interazione reale
  resta fuori dal controllo del software. \emph{Estensione:} un passo opzionale di verifica del
  resto, in cui la persona compone i pezzi ricevuti e l'app conferma che il conto torna.
  \item \textbf{Allenamento non adattivo.} Le cifre della modalit\`a allenamento sono generate
  casualmente da un insieme di tagli semplici, senza memoria dei successi e degli errori della
  persona. \emph{Estensione:} difficolt\`a progressiva guidata dallo storico, con obiettivi
  assegnabili dal tutor.
  \item \textbf{Solo valuta Euro e lingua italiana.} Tagli, fotografie e parser vocale sono
  cablati sull'Euro e sull'italiano (\texttt{EURO\_DENOMINATIONS}, \texttt{moneyImages.ts},
  numerali in \texttt{useVoiceInput}). \emph{Estensione:} astrazione della valuta e
  internazionalizzazione dei testi.
\end{itemize}

\subsection{Limiti tecnici e di piattaforma}
\begin{itemize}[leftmargin=1.4cm]
  \item \textbf{Notifiche push limitate ad Android.} \texttt{registerForPushNotifications}
  restituisce \texttt{null} su iOS: un tutor con iPhone non riceve gli avvisi, incluso l'SOS
  (l'evento resta comunque visibile in dashboard via Realtime). \emph{Estensione:} configurazione
  APNs e build di distribuzione con EAS.
  \item \textbf{OCR e voce richiedono la build nativa.} ML Kit e Speech Recognition non esistono
  in Expo Go: i moduli sono caricati con \texttt{require} difensivo e, in loro assenza, l'app
  ripiega sul solo tastierino. \emph{Estensione:} distribuzione esclusiva via dev-client o store.
  \item \textbf{Il pagamento reale richiede connettivit\`a.} Il portafoglio \`e persistito
  localmente (AsyncStorage) e consultabile offline, ma la conferma di un pagamento scrive su
  Supabase: senza rete l'operazione fallisce con un avviso. \emph{Estensione:} coda locale di
  eventi con riconciliazione al ritorno della connessione.
  \item \textbf{Costo computazionale dell'algoritmo.} La ricerca delle somme raggiungibili \`e
  pseudo-polinomiale nel numero di pezzi e nell'importo in centesimi; con inventari molto grandi
  potrebbe degradare. Nella pratica l'inventario di un esercizio \`e piccolo; un limite esplicito
  di cardinalit\`a metterebbe il vincolo nero su bianco.
  \item \textbf{Scalabilit\`a del Realtime.} La dashboard del tutor si sottoscrive agli eventi
  sulle tabelle \texttt{wallets} e \texttt{activity\_logs} filtrando lato client, e i log letti
  sono limitati alle ultime 300 righe. Con molti utenti conviene spostare il filtro lato server e
  paginare.
\end{itemize}

\subsection{Limiti di sicurezza}
\begin{itemize}[leftmargin=1.4cm]
  \item \textbf{PIN del tutor in chiaro.} Il campo \texttt{tutor\_pin} \`e memorizzato come testo
  nel profilo: \`e una barriera contro l'apertura accidentale delle impostazioni da parte dello
  studente, non un meccanismo di sicurezza forte. \emph{Estensione:} hashing del PIN o biometria
  con \texttt{expo-local-authentication}, gi\`a tra le dipendenze del progetto.
  \item \textbf{Invio push dal client.} Lo studente legge i token dei tutor associati e chiama
  direttamente l'Expo Push API: funzionale ma esposto ad abusi. \emph{Estensione:} spostare
  l'invio in una Edge Function Supabase autorizzata.
  \item \textbf{QR di associazione senza scadenza.} Il codice mostrato dallo studente \`e statico
  e l'associazione non richiede conferma reciproca. \emph{Estensione:} codici a tempo e doppio
  consenso.
\end{itemize}

\subsection{Limiti di accessibilit\`a residui}
\begin{itemize}[leftmargin=1.4cm]
  \item L'app \`e costruita per la disabilit\`a cognitiva (e, di riflesso, serve bene anche
  l'ipovisione grazie a contrasti e dimensioni), ma non affronta esplicitamente le
  \textbf{disabilit\`a motorie gravi} (switch access, scansione) n\'e \`e stata validata in modo
  sistematico con gli screen reader di sistema. \emph{Estensione:} campagna di test con
  TalkBack e VoiceOver e supporto a input alternativi.
  \item Manca una \textbf{lettura vocale in uscita} (text-to-speech) delle istruzioni: oggi
  l'output \`e visivo e aptico ma non parlato, mentre per utenti con difficolt\`a di lettura la
  voce sarebbe il canale pi\`u diretto.
  \item Soprattutto, manca una \textbf{validazione sul campo con utenti reali}: le scelte di
  accessibilit\`a rispettano gli standard, ma solo una sperimentazione con persone del target e i
  loro educatori pu\`o confermare che il carico cognitivo percepito sia davvero quello progettato.
\end{itemize}

\subsection{Sintesi}

La maggior parte dei limiti elencati discende da scelte di scope coerenti con un progetto
didattico --- valuta unica, lingua unica, push solo Android, sicurezza del PIN proporzionata alla
minaccia reale --- pi\`u che da difetti d'impianto. L'architettura scelta (logica in funzioni pure
testate, autorizzazioni dichiarate nel database con RLS, sincronizzazione a eventi) fa s\`i che le
estensioni descritte --- verifica del resto, allenamento adattivo, push iOS, coda offline, Edge
Function, internazionalizzazione --- si possano innestare in modo incrementale, senza rimettere
mano alle fondamenta.

\end{document}
